\scanpage{220}
\uuline{Autre exemple}

$X$ projectif sur $S$ loc. noeth, \unsure{,} $f_*(\underline{\mathcal O}_X)=\underline{\mathcal O}_S$

Considérons $g:\widehat C_X\longrightarrow\widehat C$ [on suppose $X$ sous la forme $\operatorname{Proj}(\underline S)$, \deleted{\unsure{$S$}} $\underline S$ engendré par $\underline S_1$]

alors $g_*(\underline{\mathcal O}_{\widehat C_X})=\underline{\mathcal O}_{\widehat C}$ [\unsure{réf|vif|vérif}]

d'où le fait que
\[
H^1(\widehat C,\underline{\mathcal O}_{\widehat C}^*)\unsure{-|\to} H^1(\widehat C_X,\mathcal O_{\widehat C_X}^*)\quad\text{est \uline{injectif}}
\]
[\unsure{réf|vérif} \unsure{§} ; sur \deleted{\uline{faisceaux inversibles}} \uline{loc. libres quelconques}]

Si \deleted{\unsure{image}} $\unsure{\xi|\zeta}$ est dans l'image, alors [désignant par $\varepsilon:X\unsure{-|\to}\widehat C_X$ la section nulle] on aura $\varepsilon^*(\xi)=0$. Or [réf] $\xi=\varphi^*(\eta)\deleted{\unsure{\mathcal O_{\widehat C_X}}}(n)$, où \deleted{$\mathcal O_{\widehat C_X}$ est le fibré \unsure{can}} $(n)$ est le fibré canonique sur $\widehat C_X$ considéré comme fibré projectif sur $X$, d'où $\varepsilon^*(\xi)=(\varphi\varepsilon)^*(\eta)\,\varepsilon^*(\mathcal O_{\widehat C_X}(n))=\eta$, on a $\eta=0$, donc $\xi=\mathcal O_{\widehat C_X}(n)$, donc $\xi$ provient de $\mathcal O_{\widehat C}(n)$. D'où
\[
\boxed{H^1(\widehat C,\underline{\mathcal O}_{\widehat C}^*)\simeq\mathbb Z,\quad\text{générateur }\operatorname{cl}(\underline{\mathcal O}_{\widehat C}(1))}
\]
% archive sheet number: 219
